\frfilename{mt246.tex} \versiondate{17.11.06} \loadeusm

\def\chaptername{Ruang fungsi} \def\sectionname{Keterintegralan seragam}

\newsection{246}

Topik berikut cukup khusus. Namun, karena alasan yang berlainan, topik ini
sangat penting baik dalam teori probabilitas maupun dalam analisis fungsional,
sehingga layak dibahas secara tuntas sejak awal.

\vleader{72pt}{246A}{Definisi} Misalkan $(X,\Sigma,\mu)$ adalah ruang ukur.

\header{246Aa}(a) Himpunan $A\subseteq{\eusm L}^1(\mu)$ {\bf
terintegralkan secara seragam} jika, untuk setiap $\epsilon>0$, terdapat
himpunan $E\in\Sigma$ dengan ukuran berhingga dan bilangan
$M\ge 0$ sedemikian sehingga

\Centerline{$\int(|f|-M\chi E)^+\le\epsilon$ untuk setiap $f\in A$.}

\header{246Ab}(b) Himpunan $A\subseteq L^1(\mu)$ {\bf
terintegralkan secara seragam} jika, untuk setiap $\epsilon>0$, terdapat
himpunan $E\in\Sigma$ dengan ukuran berhingga dan bilangan
$M\ge 0$ sedemikian sehingga

\Centerline{$\int(|u|-M\chi E^{\ssbullet})^+\le\epsilon$ untuk setiap $u\in
A$.}

\leader{246B}{Catatan} \cmmnt{{\bf (a)} Ingat rumus-rumus dari 241Ef:
$u^+=u\vee 0$, jadi $(u-v)^+=u-u\wedge v$.

\header{246Bb}{\bf (b)} Frasa `terintegralkan secara seragam' sendiri tidak
terlalu membantu. Namun, untuk setiap fungsi terintegralkan $f$, terdapat
fungsi-fungsi sederhana yang menghampiri $f$ dalam norma $\|\,\|_1$ (242M).
Modulus fungsi sederhana tersebut dibatasi oleh fungsi berbentuk $M\chi E$,
dengan $\mu E<\infty$; karena itu, setiap himpunan beranggota tunggal dalam
$\eusm L^1$ atau $L^1$ terintegralkan secara seragam. Pada himpunan fungsi
terintegralkan secara seragam yang umum, $M$ dan $E$ dapat dipilih sekaligus
untuk seluruh anggota himpunan.

\medskip

}{\bf (c)}\cmmnt{ Dari definisi, semestinya jelas bahwa }
$A\subseteq\eusm L^1$ terintegralkan secara seragam jika dan hanya
jika $\{f^{\ssbullet}:f\in A\}\subseteq L^1$ terintegralkan secara seragam.

\header{246Bd}{\bf (d)}\dvro{ Jika}{ Definisi dapat disederhanakan jika}
$\mu X<\infty$\dvro{}{ (khususnya, jika $(X,\Sigma,\mu)$ adalah ruang
probabilitas). Dalam hal ini,} himpunan $A\subseteq L^1(\mu)$
terintegralkan secara seragam jika dan hanya jika

\Centerline{$\inf_{M\ge 0}\sup_{u\in A}\int(|u|-Me)^+=0$}

\noindent jika dan hanya jika

\Centerline{$\lim_{M\to\infty}\sup_{u\in A}\int(|u|-Me)^+=0$,}

\noindent dengan $e=\chi X^{\ssbullet}\in L^1(\mu)$.
\prooflet{(Sebab jika
$\sup_{u\in A}\int(|u|-M\chi E^{\ssbullet})^+\le\epsilon$, maka
$\int(|u|-M'e)^+\le\epsilon$ untuk setiap $M'\ge M$.)} \cmmnt{Demikian pula,
}$A\subseteq\eusm L^1(\mu)$ terintegralkan secara seragam jika dan hanya jika

\Centerline{$\lim_{M\to\infty}\sup_{f\in A}\int(|f|-M\chi X)^+
=0$}

\noindent jika dan hanya jika

\Centerline{$\inf_{M\ge 0}\sup_{f\in A}\int(|f|-M\chi X)^+=0$.}

\cmmnt{\medskip

\noindent{\bf Peringatan!} Beberapa penulis menggunakan frasa `terintegralkan
secara seragam' bagi himpunan yang memenuhi syarat-syarat dalam
(d), bahkan ketika $\mu$ tidak berhingga total.}

\leader{246C}{}\cmmnt{ Kelas himpunan terintegralkan secara seragam dalam
$L^1$ atau $\eusm L^1$ memiliki berbagai sifat kestabilan berikut.

\medskip

\noindent}{\bf Proposisi} Misalkan $(X,\Sigma,\mu)$ ruang ukur dan $A$
subhimpunan $L^1(\mu)$ yang terintegralkan secara seragam.

(a) $A$ terbatas dalam norma $\|\,\|_1$.

(b) Setiap subhimpunan dari $A$ terintegralkan secara seragam.

(c) Untuk setiap $a\in\Bbb R$, $aA=\{au:u\in A\}$ terintegralkan secara seragam.

(d) Terdapat $C\supseteq A$ yang terintegralkan secara seragam sedemikian
sehingga $C$ konveks dan tertutup dalam norma $\|\,\|_1$, serta $v\in C$
setiap kali $u\in C$ dan
$|v|\le|u|$.

(e) Jika $B$ adalah himpunan terintegralkan secara seragam lainnya dalam
$L^1$, maka $A\cup B$ dan $A+B=\{u+v:u\in A,\,v\in B\}$ juga terintegralkan
secara seragam.

\proof{ Tuliskan $\Sigma^f$ untuk $\{E:E\in\Sigma,\,\mu E<\infty\}$.

\medskip

{\bf (a)} Harus ada $E\in\Sigma^f$, $M\ge 0$ sehingga $\int(|u|-M\chi
E^{\ssbullet})^+\le 1$ untuk setiap $u\in A$; sekarang

\Centerline{$\|u\|_1\le\int(|u|-M\chi E^{\ssbullet})^+
+\int M\chi
E^{\ssbullet}\le 1+M\mu E$}

\noindent untuk setiap $u\in A$, sehingga $A$ terbatas.

\medskip

{\bf (b)} Ini langsung dari definisi 246Ab.

\medskip

{\bf (c)} Untuk $\epsilon>0$, kita dapat memilih $E\in\Sigma^f$ dan $M\ge
0$ sedemikian sehingga $|a|\int(|u|-M\chi E^{\ssbullet})^+\le\epsilon$
untuk setiap $u\in A$. Karena itu,
$\int(|v|-|a|M\chi E^{\ssbullet})^+\le\epsilon$ untuk setiap $v\in aA$.

\medskip

{\bf (d)} Jika $A$ kosong, ambil $C=A$. Jika tidak, ambil

\Centerline{$C=\{v:v\in L^1,\,\int(|v|-w)^+
\le\sup_{u\in A}\int(|u|-w)^+$
untuk setiap $w\in L^1(\mu)\}$.}

\noindent Jelas bahwa $A\subseteq C$, dan $C$ memenuhi definisi 246Ab karena
$A$ memenuhinya: cukup tinjau $w$ berbentuk $M\chi E^{\ssbullet}$,
dengan $E\in\Sigma^f$ dan $M\ge 0$. Setiap fungsional

\Centerline{$v\mapsto\int(|v|-w)^+:L^1(\mu)\to\Bbb R$}

\noindent kontinu dalam norma $\|\,\|_1$ (karena operator $v\mapsto|v|$,
$v\mapsto v-w$, $v\mapsto v^+$, dan $v\mapsto\int v$ semuanya kontinu),
sehingga $C$ tertutup.
Jika $|v'|\le|v|$ dan $v\in C$, maka

\Centerline{$\int(|v'|-w)^+\le\int(|v|-w)^+
\le\sup_{u\in A}\int(|u|-w)^+$}

\noindent untuk setiap $w$, dan $v'\in C$. Jika $v=av_1+bv_2$ di mana $v_1$,
$v_2\in C$, $a\in[0,1]$ dan $b=1-a$, maka $|v|\le a|v_1|+b|v_2|$, jadi

\Centerline{$|v|-w\le(a|v_1|-aw)+(b|v_2|-bw)
\le(a|v_1|-aw)^++(b|v_2|-bw)^+$}

\noindent dan

\Centerline{$(|v|-w)^+\le a(|v_1|-w)^++b(|v_2|-w)^+$}

\noindent untuk setiap $w$; oleh karena itu,

$$\eqalign{\int(|v|-w)^+
&\le a\int(|v_1|-w)^++b\int(|v_2|-w)^+\cr
&\le(a+b)\sup_{u\in A}\int(|u|-w)^+
=\sup_{u\in A}\int(|u|-w)^+\cr}$$

\noindent untuk setiap $w$, sehingga $v\in C$.

Jadi $C$ memiliki semua sifat yang disyaratkan.

\medskip

{\bf (e)} Pertama-tama akan ditunjukkan bahwa $A\cup B$ terintegralkan secara
seragam. \Prf\ Untuk $\epsilon>0$, pilih $M_1$, $M_2\ge 0$ dan
$E_1$, $E_2\in\Sigma^f$ sedemikian sehingga

\Centerline{$\int(|u|-M_1\chi E_1^{\ssbullet})^+\le\epsilon$ untuk setiap
$u\in A$,}

\Centerline{$\int(|u|-M_2\chi E_2^{\ssbullet})^+\le\epsilon$ untuk setiap
$u\in B$.}

\noindent Tetapkan $M=\max(M_1,M_2)$ dan $E=E_1\cup E_2$; maka $\mu E<\infty$ dan

\Centerline{$\int(|u|-M\chi E^{\ssbullet})^+\le\epsilon$ untuk setiap $u\in
A\cup B$.}

\noindent Karena $\epsilon$ sembarang, $A\cup B$ terintegralkan secara
seragam. \Qed

Menurut (d), terdapat himpunan konveks terintegralkan secara seragam $C$ yang
memuat $A\cup B$. Dalam hal ini $A+B\subseteq 2C$, sehingga (b) dan (c)
menunjukkan bahwa $A+B$ juga terintegralkan secara seragam.
}%end of proof of 246C

\leader{246D}{Proposisi} Misalkan $(X,\Sigma,\mu)$ ruang probabilitas dan
$A\subseteq L^1(\mu)$ himpunan terintegralkan secara seragam. Maka terdapat
himpunan konveks, tertutup dalam norma $\|\,\|_1$, dan terintegralkan secara
seragam, yakni $C\subseteq L^1$,
sedemikian sehingga $A\subseteq C$, $w\in C$ setiap kali $v\in C$ dan
$|w|\le|v|$, dan $Pv\in C$ setiap kali $v\in C$ dan $P$ adalah operator
ekspektasi bersyarat yang terkait dengan subaljabar-sigma dari $\Sigma$.

\proof{ Tetapkan

\Centerline{$C=\{v:v\in L^1(\mu),\,
  \int(|v|-Me)^+\le\sup_{u\in
A}\int(|u|-Me)^+$ untuk setiap $M\ge 0\}$,}

\noindent dengan $e=\chi X^{\ssbullet}$ seperti biasa. Argumen dalam
bukti 246Cd menunjukkan bahwa $C\supseteq A$ terintegralkan secara seragam,
konveks, dan tertutup, serta bahwa $w\in C$ setiap kali $v\in C$ dan
$|w|\le|v|$. Adapun operator ekspektasi bersyarat, jika $v\in C$, $\Tau$
adalah subaljabar-sigma dari $\Sigma$, $P$ operator ekspektasi
bersyarat yang terkait, dan $M\ge 0$, maka

\Centerline{$|Pv|\le P|v|=P((|v|\wedge Me)+(|v|-Me)^+)
\le Me+P((|v|-Me)^+)$,}

\noindent jadi

\Centerline{$(|Pv|-Me)^+\le P((|v|-Me)^+)$}

\noindent dan

\Centerline{$\int(|Pv|-Me)^+\le\int P(|v|-Me)^+=\int(|v|-Me)^+
\le\sup_{u\in
A}\int(|u|-Me)^+$;}

\noindent karena $M$ sembarang, $Pv\in C$.
}

\cmmnt{ \leader{246E}{Catatan (a)} Tentu saja 246D dapat dirumuskan dengan
$\eusm L^1$, bukan $L^1$: jika $(X,\Sigma,\mu)$ ruang probabilitas dan
$A\subseteq\eusm L^1(\mu)$ terintegralkan secara seragam, maka terdapat
himpunan terintegralkan secara seragam $C\supseteq A$ sedemikian
sehingga (i) $af+(1-a)g\in C$ setiap kali $f$, $g\in C$ dan $a\in[0,1]$ (ii) $g\in
C$ setiap kali $f\in C$, $g\in\eusm L^0(\mu)$ dan $|g|\leae|f|$ (iii) $f\in C$
setiap kali terdapat suatu barisan $\sequencen{f_n}$ dalam $C$ sedemikian
sehingga $\lim_{n\to\infty}\int|f-f_n|=0$ (iv) $g\in C$ setiap kali terdapat
$f\in C$ sedemikian hingga $g$ merupakan ekspektasi bersyarat dari $f$
terhadap suatu subaljabar-sigma dari $\Sigma$.

\header{246Eb}{\bf (b)} Sebenarnya, 246D dapat diperluas dengan mudah;
bukti di sana sudah menunjukkan bahwa $T[C]\subseteq C$ setiap kali
$T:L^1(\mu)\to L^1(\mu)$ adalah operator linear yang melestarikan urutan,
sehingga $\|Tu\|_1\le\|u\|_1$ untuk setiap $u\in L^1(\mu)$ dan
$\|Tu\|_{\infty}\le\|u\|_{\infty}$ untuk setiap $u\in L^1(\mu)\cap
L^{\infty}(\mu)$ (246Yc). Seandainya teori operator pada ruang Riesz telah
dibahas lebih jauh, pembahasan ini pun dapat dikembangkan lebih lanjut;
misalnya, operator $T$ dalam kalimat terakhir sebenarnya tidak perlu
melestarikan urutan. Saya akan kembali ke masalah ini dalam Bab 37 pada volume
berikutnya.

\header{246Ec}{\bf (c)} Selain itu, teorema utama pada bagian berikut akan
menunjukkan bahwa untuk setiap ruang ukur $(X,\Sigma,\mu)$, $(Y,\Tau,\nu)$,
$T[A]$ terintegralkan secara seragam dalam $L^1(\nu)$ setiap kali
$A\subseteq L^1(\mu)$ terintegralkan secara seragam dan $T:L^1(\mu)\to
L^1(\nu)$ adalah operator linear kontinu (247D).
}%end of comment

\leader{246F}{}\cmmnt{ Kita memerlukan sebuah lema elementer yang belum
dinyatakan sebelumnya.

\medskip

\noindent}{\bf Lema} Misalkan $(X,\Sigma,\mu)$ ruang ukur. Maka
untuk setiap $u\in L^1(\mu)$,

\Centerline{$\|u\|_1\le 2\sup_{E\in\Sigma}|\int_Eu|$.}

\proof{ Nyatakan $u$ sebagai $f^{\ssbullet}$, dengan $f:X\to\Bbb R$ terukur.
Tetapkan $F=\{x:f(x)\ge 0\}$. Maka

\Centerline{$\|u\|_1=\int|f|=|\int_Ff|+|\int_{X\setminus F}f|
\le
2\sup_{E\in\Sigma}|\int_Ef|=2\sup_{E\in\Sigma}|\int_Eu|$.} }

\leader{246G}{}\cmmnt{ Sekarang kita sampai pada beberapa pencirian alternatif
yang mencolok bagi keterintegralan seragam.

\medskip

\noindent}{\bf Teorema} Misalkan $(X,\Sigma,\mu)$ ruang ukur sembarang dan
$A$ subhimpunan tak kosong dari $L^1(\mu)$. Maka pernyataan-pernyataan berikut
ekuivalen:

\quad(i) $A$ terintegralkan secara seragam;

\quad(ii) $\sup_{u\in A}|\int_Fu|<\infty$ untuk setiap $\mu$-atom
$F\in\Sigma$, dan untuk setiap $\epsilon>0$ terdapat $E\in\Sigma$, $\delta>0$
sedemikian rupa sehingga $\mu E<\infty$ dan $|\int_Fu|\le\epsilon$ setiap kali
$u\in A$, $F\in\Sigma$ dan $\mu(F\cap E)\le\delta$;

\quad(iii) $\sup_{u\in A}|\int_Fu|<\infty$ untuk setiap $\mu$-atom
$F\in\Sigma$, dan $\lim_{n\to\infty}\sup_{u\in A}|\int_{F_n}u|=0$ setiap kali
$\sequencen{F_n}$ barisan saling lepas dalam $\Sigma$;

\quad(iv) $\sup_{u\in A}|\int_Fu|<\infty$ untuk setiap $\mu$-atom
$F\in\Sigma$, dan $\lim_{n\to\infty}\sup_{u\in A}|\int_{F_n}u|=0$ setiap kali
$\sequencen{F_n}$ barisan tak meningkat dalam $\Sigma$ dengan irisan
kosong.

\cmmnt{\medskip

\noindent{\bf Catatan} Saya menggunakan frasa `$\mu$-atom' untuk menekankan
bahwa saya maksud atom dalam arti ruang ukur (211I).}

\proof{{\bf (a)(i)$\Rightarrow$(iv)} Misalkan $A$ terintegralkan secara
seragam. Jelas bahwa jika $F\in\Sigma$ adalah sebuah $\mu$-atom,

\Centerline{$\sup_{u\in A}|\int_Fu|
\le\sup_{u\in A}\|u\|_1<\infty$,}

\noindent menurut 246Ca. Misalkan $\sequencen{F_n}$ barisan tak
meningkat dalam $\Sigma$ dengan irisan kosong, dan misalkan
$\epsilon>0$. Ambil $E\in\Sigma$, $M\ge 0$ sedemikian rupa sehingga $\mu
E<\infty$ dan $\int(|u|-M\chi E^{\ssbullet})^+\le\bover12\epsilon$ setiap kali
$u\in A$. Untuk semua $n$ yang cukup besar berlaku $M\mu(F_n\cap
E)\le\bover12\epsilon$, sehingga

\Centerline{$|\int_{F_n}u|\le\int_{F_n}|u|
\le\int(|u|-M\chi
E^{\ssbullet})^++\int_{F_n}M\chi E^{\ssbullet}
\le\Bover{\epsilon}2+M\mu(F_n\cap E)\le\epsilon$}

\noindent untuk setiap $u\in A$. Karena $\epsilon$ sembarang,
$\lim_{n\to\infty}\sup_{u\in A}|\int_{F_n}u|=0$, dan (iv) benar.

\medskip

{\bf (b)(iv)$\Rightarrow$(iii)} Misalkan (iv) benar. Maka tentu saja
$\sup_{u\in A}|\int_Fu|<\infty$ untuk setiap $\mu$-atom $F\in\Sigma$. \Quer\
Andaikan, jika mungkin, bahwa $\sequencen{F_n}$ barisan saling lepas dalam
$\Sigma$ sehingga \discrcenter{468pt}{$\epsilon
=\limsup_{n\to\infty}\sup_{u\in A}\min(1,\bover13|\int_{F_n}u|)
>0$. }
Tetapkan $H_n=\bigcup_{i\ge n}F_i$ untuk setiap $n$, sehingga
$\sequencen{H_n}$ tak meningkat dan beririsan kosong, serta
$\int_{H_n}u\to 0$ sebagai $n\to\infty$ untuk setiap $u\in L^1(\mu)$. Pilih
$\sequence{i}{n_i}$, $\sequence{i}{m_i}$, $\sequence{i}{u_i}$ secara induktif
sebagai berikut. Tetapkan $n_0=0$. Setelah $n_i\in\Bbb N$ dipilih, ambil
$m_i\ge n_i$,
$u_i\in A$ sedemikian sehingga $|\int_{F_{m_i}}u_i|\ge 2\epsilon$. Ambil
$n_{i+1}>m_i$ sedemikian sehingga $\int_{H_{n_{i+1}}}|u_i|\le\epsilon$.
Lanjutkan.

Tetapkan $G_k=\bigcup_{i\ge k}F_{m_i}$ untuk setiap $k$. Maka
$\sequence{k}{G_k}$ barisan tak meningkat dalam $\Sigma$ dengan
irisan kosong. Tetapi $F_{m_i}\subseteq G_i\subseteq F_{m_i}\cup
H_{n_{i+1}}$,
jadi

\Centerline{$|\int_{G_i}u_i|
\ge|\int_{F_{m_i}}u_i|-|\int_{G_i\setminus
F_{m_i}}u_i|
\ge 2\epsilon-\int_{H_{n_{i+1}}}|u_i|
\ge\epsilon$}

\noindent untuk setiap $i$; ini bertentangan dengan hipotesis (iv). \Bang

Ini berarti bahwa $\lim_{n\to\infty}\sup_{u\in A}|\int_{F_n}u|$ harus nol, dan
(iii) berlaku.

\medskip

{\bf (c)(iii)$\Rightarrow$(ii)} Kita tetap mempunyai $\sup_{u\in
A}|\int_Fu|<\infty$ untuk setiap $\mu$-atom $F$. \Quer\ Andaikan terdapat
$\epsilon>0$ sedemikian sehingga untuk setiap
himpunan terukur $E$ berukuran berhingga dan setiap $\delta>0$ terdapat
$u\in A$,
$F\in\Sigma$ sedemikian rupa sehingga $\mu(F\cap E)\le\delta$ dan
$|\int_Fu|\ge\epsilon$. Pilih barisan $\sequencen{E_n}$ dari
himpunan-himpunan berukuran berhingga, barisan $\sequencen{G_n}$ dalam
$\Sigma$, barisan $\sequencen{\delta_n}$ bilangan real positif tegas, dan
barisan $\sequencen{u_n}$ dalam $A$ sebagai berikut.
Setelah $u_k$, $E_k$, $\delta_k$ dipilih untuk $k<n$, pilih $u_n\in A$ dan
$G_n\in\Sigma$ sedemikian rupa sehingga $\mu(G_n\cap\bigcup_{k<n}E_k)\le
2^{-n}\min(\{1\}\cup\{\delta_k:k<n\})$ dan $|\int_{G_n}u_n|\ge\epsilon$;
lalu pilih himpunan $E_n$ berukuran berhingga dan $\delta_n>0$
sedemikian sehingga $\int_F|u_n|\le\bover12\epsilon$ setiap kali $F\in\Sigma$
dan $\mu(F\cap E_n)\le\delta_n$ (lihat 225A). Lanjutkan.

Setelah induksi selesai, tetapkan $F_n=E_n\cap
G_n\setminus\bigcup_{k>n}G_k$ untuk setiap $n$; maka $\sequencen{F_n}$
merupakan barisan saling lepas dalam $\Sigma$. Berdasarkan pemilihan $G_k$,

\Centerline{$\mu(E_n\cap\bigcup_{k>n}G_k)
\le\sum_{k=n+1}^{\infty}2^{-k}\delta_n\le\delta_n$,}

\noindent sehingga $\mu(E_n\cap(G_n\setminus F_n))\le\delta_n$ dan
$\int_{G_n\setminus F_n}|u_n|\le \bover12\epsilon$. Ini berarti bahwa
$|\int_{F_n}u_n|\ge|\int_{G_n}u_n|-\bover12\epsilon\ge\bover12\epsilon$.
Tetapi ini bertentangan dengan hipotesis (iii). \Bang

\medskip

{\bf (d)(ii)$\Rightarrow$(i)}\grheada\ Misalkan (ii) berlaku dan $\epsilon>0$.
Maka terdapat $E\in\Sigma$, $\delta>0$ sedemikian sehingga $\mu E<\infty$ dan
$|\int_Fu|\le\epsilon$ setiap kali $u\in A$, $F\in\Sigma$ dan $\mu(F\cap
E)\le\delta$. Sekarang $\sup_{u\in A}\int_E|u|<\infty$. \Prf\
 Tuliskan $\Cal I$ sebagai keluarga semua $F\in\Sigma$ sedemikian sehingga
$F\subseteq E$ dan $\sup_{u\in A}\int_F|u|$ berhingga. Jika $F\subseteq E$
merupakan atom bagi $\mu$, maka
$\sup_{u\in A}\int_F|u|=\sup_{u\in A}|\int_Fu|<\infty$,
sehingga $F\in\Cal I$. (Intinya adalah bahwa jika $f:X\to\Bbb R$ adalah fungsi
terukur sedemikian sehingga $f^{\ssbullet}=u$, maka salah satu dari
$F'=\{x:x\in F,\,f(x)\ge 0\}$, $F''=\{x:x\in F,\,f(x)<0\}$ harus terabaikan,
sehingga $\int_F|u|$ adalah $\int_{F'}u=\int_Fu$ atau
$-\int_{F''}u=-\int_Fu$.) Jika $F\in\Sigma$, $F\subseteq E$ dan $\mu
F\le\delta$ maka

\Centerline{$\sup_{u\in A}\int_F|u|
\le 2\sup_{u\in A,G\in\Sigma,G\subseteq
F}|\int_Gu|
\le 2\epsilon$}

\noindent (menurut 246F), sehingga $F\in\Cal I$. Selanjutnya, jika
$F$, $G\in\Cal I$, maka $\sup_{u\in A}\int_{F\cup G}|u|
\le\sup_{u\in A}\int_F|u|+\sup_{u\in
A}\int_G|u|$ berhingga, jadi $F\cup G\in \Cal I$. Akhirnya, jika
$\sequencen{F_n}$ sembarang barisan dalam $\Cal I$, dan
$F=\bigcup_{n\in\Bbb N}F_n$, terdapat $n\in\Bbb N$ sehingga
$\mu(F\setminus\bigcup_{i\le n}F_i)\le\delta$; sekarang $\bigcup_{i\le n}F_i$
dan $F\setminus\bigcup_{i\le n}F_i$ keduanya termasuk dalam $\Cal I$, jadi
$F\in\Cal I$.

Menurut 215Ab, terdapat $F\in\Cal I$ sedemikian sehingga $H\setminus F$
terabaikan untuk setiap $H\in\Cal I$. Sekarang perhatikan bahwa
$E\setminus F$ tidak dapat memuat anggota $\Cal I$ yang tidak terabaikan;
khususnya, ia tidak dapat memuat atom ataupun himpunan tidak terabaikan yang
berukuran kurang dari $\delta$. Dengan demikian, ukuran pada subruang
$E\setminus F$ tak beratom, berhingga total, dan tidak mempunyai
himpunan terukur yang tidak terabaikan dengan ukuran
kurang dari $\delta$; berdasarkan 215D, $\mu(E\setminus F)=0$ dan baik
$E\setminus F$ maupun $E$ termasuk dalam $\Cal I$, sebagaimana diperlukan.\ \Qed

Karena $\int_{X\setminus E}|u|\le 2\epsilon$ untuk setiap $u\in A$,
$\gamma=\sup_{u\in A}\int|u|$ berhingga.

\medskip

\quad\grheadb\
Tetapkan $M=\gamma/\delta$. Jika $u\in A$, nyatakan $u$ sebagai
$f^{\ssbullet}$, dengan $f:X\to\Bbb R$ terukur, dan pertimbangkan

\Centerline{$F=\{x:f(x)\ge M\chi E(x)\}$.}

\noindent Maka

\Centerline{$M\mu(F\cap E)\le\int_Ff=\int_Fu\le\gamma$,}

\noindent sehingga $\mu(F\cap E)\le\gamma/M=\delta$. Oleh karena itu
$\int_Fu\le\epsilon$. Demikian pula, $\int_{F'}(-u)\le\epsilon$, dengan
$F'=\{x:-f(x)\ge M\chi E(x)\}$. Ini berarti bahwa

\Centerline{$\int(|u|-M\chi E^{\ssbullet})^+
=\int(|f|-M\chi
E)^+\le\int_{F\cup F'}|f|
=\int_{F\cup F'}|u|\le 2\epsilon$}

\noindent untuk setiap $u\in A$. Karena $\epsilon$ sembarang, $A$
terintegralkan secara seragam.
}%end of proof of 246G

\cmmnt{ \leader{246H}{Catatan (a)} Seperti (i), syarat (ii)--(iv) dalam
teorema ini tentu memiliki padanan langsung bagi anggota $\eusm L^1$.
Dengan demikian, himpunan tak kosong $A\subseteq\eusm L^1$
terintegralkan secara seragam jika dan hanya jika $\sup_{f\in A}|\int_Ff|$
berhingga untuk setiap atom $F\in\Sigma$ dan

\inset{{\it baik } untuk setiap $\epsilon>0$ terdapat
$E\in\Sigma$, $\delta>0$ sehingga $\mu E<\infty$ dan $|\int_Ff|\le\epsilon$
setiap kali $f\in A$, $F\in\Sigma$ dan $\mu(F\cap
E)\le\delta$ }

\inset{{\it atau } $\lim_{n\to\infty}\sup_{f\in A}|\int_{F_n}f|=0$ untuk
setiap barisan saling lepas $\sequencen{F_n}$ dalam $\Sigma$}

\inset{{\it atau } $\lim_{n\to\infty}\sup_{f\in A}|\int_{F_n}f|=0$ untuk
setiap barisan tak meningkat $\sequencen{F_n}$ dalam $\Sigma$ dengan
irisan kosong.}

\header{246Hb}{\bf (b)} Masih ada tak terhitung banyaknya ungkapan ekuivalen
yang mencirikan keterintegralan seragam; setiap penulis memiliki pilihannya
sendiri. Banyak di antaranya merupakan variasi atas (i)--(iv) dalam teorema
ini, seperti pada 246I dan 246Yd--246Yf. Untuk syarat yang sama sekali berbeda,
lihat Teorema 247C.
}%end of comment

\leader{246I}{Akibat} Misalkan $(X,\Sigma,\mu)$ ruang probabilitas. Untuk
$f\in\eusm L^0(\mu)$, $M\ge 0$, tetapkan
$F(f,M)=\{x:x\in\dom f,\,|f(x)|\ge M\}$. Maka himpunan tak kosong
$A\subseteq\eusm L^1(\mu)$
terintegralkan secara seragam jika dan hanya jika
\discrcenter{468pt}{$\lim_{M\to\infty}\sup_{f\in A}\int_{F(f,M)}|f|=0$.}

\proof{{\bf (a)} Jika $A$ memenuhi syarat tersebut, maka

\Centerline{$\inf_{M\ge 0}\sup_{f\in A}\int(|f|-M\chi X)^+
\le\inf_{M\ge
0}\sup_{f\in A}\int_{F(f,M)}|f|=0$,}

\noindent sehingga $A$ terintegralkan secara seragam.

\medskip

{\bf (b)} Jika $A$ terintegralkan secara seragam dan $\epsilon>0$, terdapat
$M_0\ge 0$ sedemikian sehingga $\int(|f|-M_0\chi X)^+\le\epsilon$ untuk setiap
$f\in
A$; selain itu, $\gamma=\sup_{f\in A}\int|f|$ berhingga (246Ca). Ambil sembarang
$M\ge M_0\max(1,(1+\gamma)/\epsilon)$. Jika $f\in A$, maka

\Centerline{$|f|\times\chi F(f,M)
\le (|f|-M_0\chi X)^++M_0\chi F(f,M)
\le
(|f|-M_0\chi X)^++\Bover{\epsilon}{\gamma+1}|f|$}

\noindent di seluruh $\dom f$, sehingga

\Centerline{$\int_{F(f,M)}|f|
\le\int(|f|-M_0\chi
X)^++\Bover{\epsilon}{\gamma+1}\int|f|
\le 2\epsilon$.}

\noindent Karena $\epsilon$ sembarang,
$\lim_{M\to\infty}\sup_{f\in
A}\int_{F(f,M)}|f|=0$.
}%end of proof of 246I

\leader{246J}{}\cmmnt{ Langkah berikutnya ialah memaparkan beberapa kaitan
penting antara keterintegralan seragam dan topologi konvergensi dalam
ukuran yang dibahas dalam bagian sebelumnya.

\medskip

\noindent}{\bf Teorema} Misalkan $(X,\Sigma,\mu)$ ruang ukur.

\ifdim\pagewidth>467pt\fontdimen4\tenrm=1.5pt\fi (a) Jika $\sequencen{f_n}$
adalah barisan fungsi bernilai real pada $X$ yang terintegralkan secara
seragam dan $f(x)=\lim_{n\to\infty}f_n(x)$ untuk hampir setiap $x\in X$, maka $f$
terintegralkan dan $\lim_{n\to\infty}\int|f_n-f|=0$; akibatnya $\int
f=\lim_{n\to\infty}\int f_n$. \fontdimen4\tenrm=1.11pt

(b) Jika $A\subseteq L^1=L^1(\mu)$ terintegralkan secara seragam, maka
topologi norma pada $L^1$ dan topologi konvergensi dalam ukuran pada
$L^0=L^0(\mu)$ berimpit pada $A$.

(c) Untuk setiap $u\in L^1$ dan setiap barisan $\sequencen{u_n}$ dalam $L^1$,
pernyataan-pernyataan berikut ekuivalen:

\qquad(i) $u=\lim_{n\to\infty}u_n$ dalam norma $\|\,\|_1$;

\qquad(ii) $\{u_n:n\in\Bbb N\}$ terintegralkan secara seragam dan
$\sequencen{u_n}$ konvergen ke $u$ dalam ukuran.

(d) Jika $(X,\Sigma,\mu)$ semihingga dan $A\subseteq L^1$ terintegralkan secara
seragam, maka penutupan $\overline{A}$ dari $A$ dalam $L^0$, untuk topologi
konvergensi dalam ukuran, tetap merupakan subhimpunan $L^1$ yang terintegralkan
secara seragam.

\proof{{\bf (a)} Perhatikan dahulu bahwa karena $\sup_{n\in\Bbb
N}\int|f_n|<\infty$ (246Ca) dan $|f|=\liminf_{n\to\infty}|f_n|$, Lemma Fatou
memastikan bahwa $|f|$ terintegralkan, dengan
$\int|f|\le\limsup_{n\to\infty}\int|f_n|$. Dengan segera diperoleh bahwa
$\{f_n-f:n\in\Bbb N\}$ terintegralkan secara
seragam, karena merupakan jumlah dua himpunan terintegralkan secara seragam
(246Cc, 246Ce).

Untuk $\epsilon>0$, terdapat $M\ge 0$, $E\in\Sigma$ sehingga $\mu
E<\infty$ dan
$\int(|f_n-f|-M\chi E)^+\le\epsilon$ untuk setiap $n\in\Bbb
N$. Selain itu,
$|f_n-f|\wedge M\chi E\to 0$ hampir di mana-mana, sehingga

$$\eqalign{\limsup_{n\to\infty}\int|f_n-f|
&\le\limsup_{n\to\infty}\int(|f_n-f|-M\chi E)^+\cr
&\qquad\qquad+\limsup_{n\to\infty}\int|f_n-f|\wedge M\chi E\cr
&\le\epsilon,\cr}$$

\noindent menurut Teorema Konvergensi Terdominasi Lebesgue. Karena $\epsilon$
sembarang, $\lim_{n\to\infty}\int|f_n-f|=0$ dan $\lim_{n\to\infty}\int
f_n-f=0$.

\medskip

{\bf (b)} Misalkan $\frak T_A$ dan $\frak S_A$ berturut-turut adalah topologi
pada $A$ yang diinduksi oleh topologi norma pada $L^1$ dan topologi
konvergensi dalam ukuran pada $L^0$.

\medskip

\quad{\bf (i)} Untuk $\epsilon>0$, pilih $F\in\Sigma$ dan $M\ge 0$
sedemikian sehingga $\mu F<\infty$ dan $\int(|v|-M\chi
F^{\ssbullet})^+\le\epsilon$ untuk setiap $v\in A$, dan pertimbangkan
$\bar\rho_F$, yang didefinisikan seperti di 245A. Kemudian untuk setiap $f$,
$g\in\eusm L^0$,

\Centerline{$|f-g|\le(|f|-M\chi F)^+
+(|g|-M\chi F)^++M(|f-g|\wedge \chi F)$}

\noindent di seluruh $\dom f\cap \dom g$, sehingga

\Centerline{$|u-v|\le(|u|-M\chi F^{\ssbullet})^+
+(|v|-M\chi
F^{\ssbullet})^++M(|u-v|\wedge\chi F^{\ssbullet})$}

\noindent untuk semua $u$, $v\in L^0$. Akibatnya

\Centerline{$\|u-v\|_1\le 2\epsilon+M\bar\rho_F(u,v)$}

\noindent untuk semua $u$, $v\in A$.

Artinya, untuk $\epsilon>0$, kita dapat memilih $F$ dan $M$ sehingga, bagi
$u$, $v\in A$,

\Centerline{$\bar\rho_F(u,v)
\le\Bover{\epsilon}{1+M}\Longrightarrow\|u-v\|_1\le 3\epsilon$.}

\noindent Maka setiap subhimpunan dari $A$ yang terbuka dalam $\frak T_A$
juga terbuka dalam $\frak S_A$ (2A3Ib).

\medskip

\quad{\bf (ii)} Sebaliknya, kita mempunyai $\bar\rho_F(u,v)\le\|u-v\|_1$
untuk setiap $u\in L^1$ dan setiap himpunan $F$ berukuran berhingga,
sehingga setiap subhimpunan dari $A$ yang terbuka dalam $\frak S_A$
terbuka dalam $\frak T_A$.

\medskip

{\bf (c)} Jika $\sequencen{u_n}\to u$ dalam norma $\|\,\|_1$, maka
$A=\{u_n:n\in\Bbb N\}$ terintegralkan secara seragam. \Prf\ Untuk
$\epsilon>0$,
pilih $m$ sedemikian rupa sehingga $\|u_n-u\|_1\le\epsilon$ setiap kali
$n\ge m$. Tetapkan $v=|u|+\sum_{i\le m}|u_i|\in L^1$, dan pilih $M\ge 0$,
$E\in\Sigma$ sedemikian sehingga $\mu E$ berhingga dan
$\int(v-M\chi E^{\ssbullet})^+\le\epsilon$. Maka, untuk $w\in A$,

\Centerline{$(|w|-M\chi E^{\ssbullet})^+
\le(|w|-v)^++(v-M\chi
E^{\ssbullet})^+$,}

\noindent jadi

\Centerline{$\int(|w|-M\chi E^{\ssbullet})^+
\le\|(|w|-v)^+\|_1+\int(v-M\chi E^{\ssbullet})^+
\le 2\epsilon$. \Qed}

Jadi, di bawah salah satu hipotesis, $\{u_n:n\in\Bbb N\}$ dan
$A=\{u\}\cup\{u_n:n\in\Bbb N\}$ terintegralkan secara seragam. Karena itu,
kedua topologi berimpit pada $A$ (menurut (b)), dan
$\sequencen{u_n}$ konvergen
ke $u$ dalam satu topologi jika dan hanya jika ia konvergen ke $u$ dalam
topologi yang lain.

\medskip

{\bf (d)} Karena $A$ terbatas dalam norma $\|\,\|_1$ (246Ca) dan $\mu$
semihingga, $\overline{A}\subseteq L^1$ (245J(b-i)). Untuk
$\epsilon>0$, pilih $M\ge 0$, $E\in\Sigma$ sedemikian sehingga $\mu
E<\infty$ dan $\int(|u|-M\chi E^{\ssbullet})^+\le\epsilon$ untuk setiap $u\in
A$.   Sekarang peta $u\mapsto|u|$, $u\mapsto u-M\chi E^{\ssbullet}$, $u\mapsto
u^+:L^0\to L^0$ semuanya kontinu untuk topologi konvergensi dalam ukuran
(245D), sementara $\{u:\|u\|_1\le\epsilon\}$ tertutup untuk topologi yang sama
(245J lagi), sehingga $\{u:u\in L^0,\,\int(|u|-M\chi
E^{\ssbullet})^+\le\epsilon\}$ tertutup dan harus mencakup $\overline{A}$.
Dengan demikian $\int(|u|-M\chi E^{\ssbullet})^+\le\epsilon$ untuk setiap
$u\in\overline{A}$. Karena $\epsilon$ sembarang, $\overline{A}$
terintegralkan secara seragam.
}%end of proof of 246J

\leader{246K}{$\eusm L^1$ dan $L^1$ kompleks}\dvro{ Untuk}{ Definisi dan
teorema di atas dapat diterapkan tanpa kesulitan pada ruang (kelas ekuivalensi
dari) fungsi bernilai kompleks, dengan satu perubahan saja: konstanta pada
padanan kompleks 246F harus diubah. Mudah dilihat bahwa, untuk} $u\in
L^1_{\Bbb C}(\mu)$, \dvro{$\|u\|_1\le 4\sup_{F\in\Sigma}|\int_Fu|$.}{

$$\eqalign{\|u\|_1
&\le \|\Real(u)\|_1+\|\Imag(u)\|_1\cr
&\le
2\sup_{F\in\Sigma}|\int_F\Real(u)|
   +2\sup_{F\in\Sigma}|\int_F\Imag(u)|
\le
4\sup_{F\in\Sigma}|\int_Fu|.\cr}$$ }

\cmmnt{\noindent (Sebenarnya, $\|u\|_1\le\pi\sup_{F\in\Sigma}|\int_Fu|$; lihat
246Yl dan 252Yt.) Akibatnya, beberapa argumen dalam 246G perlu ditulis dengan
konstanta berbeda, tetapi hasil-hasil sebagaimana dinyatakan tetap berlaku.
}%end of comment

\exercises{ \leader{246X}{Latihan dasar (a)}
%\spheader 246Xa
Misalkan $(X,\Sigma,\mu)$ ruang ukur dan $A$ subhimpunan
$L^1=L^1(\mu)$. Tunjukkan bahwa pernyataan-pernyataan berikut ekuivalen: (i) $A$
terintegralkan secara seragam; (ii) untuk setiap $\epsilon>0$ terdapat $w\ge
0$ di $L^1$ sedemikian sehingga $\int(|u|-w)^+\le\epsilon$ untuk setiap $u\in
A$; (iii) $\sequencen{(|u_{n+1}|-\sup_{i\le n}|u_i|)^+}\to 0$ dalam $L^1$
untuk setiap barisan $\sequencen{u_n}$ dalam $A$. \Hint{ untuk (ii) $\Rightarrow$
(iii), tetapkan $v_n=\sup_{i\le n}|u_i|$ dan perhatikan bahwa
$\sequencen{v_n\wedge w}$ konvergen dalam $L^1$ untuk setiap $w\ge 0$. }

\sqheader 246Xb Misalkan $(X,\Sigma,\mu)$ ruang ukur berhingga total.
Tunjukkan bahwa untuk setiap $p>1$ dan $M\ge 0$, himpunan
$\{f:f\in{\eusm L}^p(\mu),\,\|f\|_p\le M\}$ terintegralkan secara
seragam. \Hint{$\int(|f|-M\chi X)^+\le M^{1-p}\int|f|^p$.}
%246B used in \zz5R query

\sqheader 246Xc Misalkan $\mu$ ukuran hitung pada $\Bbb N$. Tunjukkan bahwa
suatu himpunan $A\subseteq\eusm L^1(\mu)=\ell^1$ terintegralkan secara
seragam jika dan hanya jika (i) $\sup_{f\in A}|f(n)|<\infty$ untuk
setiap $n\in \Bbb N$ (ii) untuk setiap $\epsilon>0$ terdapat $m\in\Bbb N$ sehingga
$\sum_{n=m}^{\infty}|f(n)|\le\epsilon$ untuk setiap $f\in A$.
%246G

\spheader 246Xd Misalkan $X$ suatu himpunan dan $\mu$ ukuran hitung pada $X$.
Tunjukkan bahwa himpunan $A\subseteq\eusm L^1(\mu)=\ell^1(X)$
terintegralkan secara seragam jika dan hanya
jika (i) $\sup_{f\in A}|f(x)|<\infty$ untuk setiap $x\in X$ (ii) untuk setiap
$\epsilon>0$ terdapat himpunan berhingga $I\subseteq
X$ sehingga $\sum_{x\in
X\setminus I}|f(x)|\le\epsilon$ untuk setiap $f\in
A$. Tunjukkan bahwa dalam hal ini $A$ relatif kompak dalam topologi norma pada
$\ell^1(X)$.
%246G

\spheader 246Xe Misalkan $(X,\Sigma,\mu)$ ruang ukur, $\delta>0$, dan
$\Cal I\subseteq\Sigma$ suatu keluarga sedemikian sehingga (i) setiap
atom termasuk dalam $\Cal I$ (ii) $E\in\Cal I$ setiap kali $E\in\Sigma$ dan
$\mu E\le\delta$ (iii) $E\cup F\in\Cal I$ setiap kali $E$, $F\in\Cal I$ dan
$E\cap
F=\emptyset$. Tunjukkan bahwa setiap himpunan berukuran berhingga
termasuk dalam $\Cal I$.
%246G

\spheader 246Xf Misalkan $(X,\Sigma,\mu)$ dan $(Y,\Tau,\nu)$ ruang ukur,
serta $\phi:X\to Y$ fungsi \imp\ . Tunjukkan bahwa suatu himpunan
$A\subseteq\eusm L^1(\nu)$ terintegralkan secara seragam jika dan hanya
jika $\{g\phi:g\in A\}$ terintegralkan secara seragam dalam
$\eusm L^1(\mu)$. \Hint{
gunakan 246G untuk `jika', 246A untuk `hanya jika'.}
%246G

\sqheader 246Xg Misalkan $(X,\Sigma,\mu)$ ruang ukur dan
$p\in\coint{1,\infty}$. Misalkan $\sequencen{f_n}$ suatu barisan dalam
$\eusm L^p=\eusm L^p(\mu)$ sedemikian sehingga $\{|f_n|^p:n\in\Bbb N\}$
terintegralkan secara seragam dan $f_n\to f$ hampir di mana-mana. Tunjukkan
bahwa $f\in\eusm L^p$ dan $\lim_{n\to\infty}\int|f_n-f|^p=0$.
%246J

\spheader 246Xh Misalkan $(X,\Sigma,\mu)$ ruang ukur semihingga dan
$p\in\coint{1,\infty}$. Misalkan $\sequencen{u_n}$ suatu barisan dalam
$L^p=L^p(\mu)$ dan $u\in L^0(\mu)$. Tunjukkan bahwa pernyataan-pernyataan
berikut ekuivalen: (i) $u\in L^p$ dan $\sequencen{u_n}$ konvergen ke $u$
dalam norma $\|\,\|_p$
(ii) $\sequencen{u_n}$ konvergen dalam ukuran ke $u$ dan $\{|u_n|^p:n\in\Bbb
N\}$ terintegralkan secara seragam. \Hint{245Xl.}
%246J

\spheader 246Xi Misalkan $(X,\Sigma,\mu)$ ruang ukur berhingga total, dan
$1\le p<r\le\infty$. Misalkan $\sequencen{u_n}$ barisan yang terbatas dalam
norma $\|\,\|_r$ dalam $L^r(\mu)$ dan konvergen dalam ukuran ke $u\in
L^0(\mu)$. Tunjukkan bahwa $\sequencen{u_n}$ konvergen ke $u$ dalam norma
$\|\,\|_p$. \Hint{ tunjukkan bahwa $\{|u_n|^p:n\in\Bbb N\}$
terintegralkan secara seragam. }
%246J, 246Xh

\leader{246Y}{Latihan lebih lanjut (a)}
%\spheader 246Ya
\ifdim\pagewidth>467pt\fontdimen4\tenrm=1.5pt\fi Misalkan $(X,\Sigma,\mu)$
ruang ukur berhingga total. Tunjukkan bahwa $A\subseteq\eusm
L^1(\mu)$ terintegralkan secara seragam jika dan hanya jika terdapat fungsi
konveks $\phi:\coint{0,\infty}\to\Bbb R$ sedemikian sehingga
$\lim_{a\to\infty}\phi(a)/a=\infty$ dan $\sup_{f\in A}\int\phi(|f|)<\infty$.
\fontdimen4\tenrm=1.11pt
%246B

\spheader 246Yb Untuk setiap ruang metrik $(Z,\rho)$, misalkan $\Cal C_Z$
keluarga subhimpunan tertutup dari $Z$. Untuk $F$,
$F'\in\Cal C_Z\setminus\{\emptyset\}$ tetapkan
$\tilde\rho(F,F')=\min(1,\max(\sup_{z\in F}\inf_{z'\in F'}\rho(z,z'),
\sup_{z'\in F'}\inf_{z\in F}\rho(z,z')))$. Tunjukkan bahwa $\tilde\rho$
merupakan metrik pada $\Cal C_Z\setminus\{\emptyset\}$ (inilah {\bf metrik
Hausdorff}). Tunjukkan bahwa jika $(Z,\rho)$ lengkap, maka keluarga
$\Cal K_Z\setminus\{\emptyset\}$ dari subhimpunan kompak tak kosong dalam $Z$
tertutup dalam metrik $\tilde\rho$. Sekarang misalkan $(X,\Sigma,\mu)$ ruang ukur
sembarang dan ambil $Z=L^1=L^1(\mu)$,
$\rho(z,z')=\|z-z'\|_1$ untuk $z$, $z'\in Z$. Tunjukkan bahwa keluarga
subhimpunan tertutup tak kosong dalam $L^1$ yang terintegralkan secara seragam
merupakan subhimpunan tertutup dari
$\Cal C_Z\setminus\{\emptyset\}$ yang memuat
$\Cal K_Z\setminus\{\emptyset\}$.
%246B

\spheader 246Yc Misalkan $(X,\Sigma,\mu)$ ruang ukur berhingga total dan
$A\subseteq L^1(\mu)$ himpunan terintegralkan secara seragam. Tunjukkan bahwa
terdapat himpunan terintegralkan secara seragam $C\supseteq A$ sedemikian
sehingga (i) $C$ konveks dan
tertutup dalam $L^0(\mu)$ untuk topologi konvergensi dalam ukuran (ii) jika
$u\in C$ dan $|v|\le|u|$ maka $v\in C$ (iii) jika $T$ termasuk dalam himpunan
$\Cal T^+$ dari operator dari $L^1(\mu)=M^{1,\infty}(\mu)$ ke dirinya sendiri,
seperti yang dijelaskan dalam 244Xm, maka $T[C]\subseteq C$.
%246C

%\ifdim\pagewidth>467pt\fontdimen4\tenrm=1.5pt\fi
\spheader 246Yd Misalkan $\mu$ ukuran Lebesgue pada $\Bbb R$. Tunjukkan bahwa
suatu himpunan $A\subseteq\eusm L^1(\mu)$ terintegralkan secara seragam
jika dan hanya jika $\lim_{n\to\infty}\int_{F_n}f_n
\ifdim\pagewidth>467pt\penalty-100\fi
=0$ untuk setiap barisan saling lepas $\sequencen{F_n}$ dari himpunan kompak
dalam $\Bbb R$ dan setiap barisan $\sequencen{f_n}$ dalam $A$.
%246G
%\fontdimen4\tenrm=1.11pt

\spheader 246Ye Misalkan $\mu$ ukuran Lebesgue pada $\Bbb R$. Tunjukkan
bahwa himpunan $A\subseteq\eusm L^1(\mu)$ terintegralkan
secara seragam jika dan hanya jika $\lim_{n\to\infty}\int_{G_n}f_n
\ifdim\pagewidth>467pt\penalty-100\fi
=0$ untuk setiap barisan saling lepas $\sequencen{G_n}$ dari himpunan terbuka
dalam $\Bbb R$ dan setiap barisan $\sequencen{f_n}$ dalam $A$.
%246G, 246Yd

\spheader 246Yf Ulangi 246Yd dan 246Ye untuk ukuran Lebesgue pada subhimpunan
sembarang dari $\BbbR^r$.
%246G, 246Yd, 246Ye

\spheader 246Yg Misalkan $X$ suatu himpunan dan $\Sigma$ aljabar-sigma dari
subhimpunan-subhimpunan $X$. Misalkan $\sequencen{\nu_n}$ barisan fungsional
aditif terhitung pada $\Sigma$ sedemikian sehingga
$\nu E=\lim_{n\to\infty}\nu_nE$ terdefinisi untuk setiap $E\in\Sigma$.
Tunjukkan bahwa $\lim_{n\to\infty}\nu_nF_n=0$ setiap kali $\sequencen{F_n}$
barisan saling lepas dalam $\Sigma$. \Hint{Andaikan sebaliknya. Dengan
mengambil subbarisan yang sesuai, cukup tinjau kasus ketika $|\nu_nF_i-\nu
F_i|\le 2^{-n}\epsilon$ untuk $i<n$, $|\nu_nF_n|\ge 3\epsilon$, $|\nu_nF_i|\le
2^{-i}\epsilon$ untuk $i>n$. Tetapkan $F=\bigcup_{i\in\Bbb N}F_{2i+1}$ dan
tunjukkan bahwa $|\nu_{2n+1}F-\nu_{2n}F|\ge\epsilon$ untuk setiap $n$.}
%246G %231Xf

\spheader 246Yh Misalkan $(X,\Sigma,\mu)$ adalah ruang ukur dan
$\sequencen{u_n}$ suatu barisan dalam $L^1=L^1(\mu)$ sehingga
$\lim_{n\to\infty}\int_Fu_n$ terdefinisi untuk setiap $F\in\Sigma$.
Tunjukkan bahwa $\{u_n:n\in\Bbb N\}$ terintegralkan secara seragam.
\Hint{Andaikan tidak. Maka terdapat barisan saling lepas
$\sequencen{F_n}$ dalam $\Sigma$ dan subbarisan $\sequencen{u'_n}$ dari
$\sequencen{u_n}$ sehingga $\inf_{n\in\Bbb N}|\int_{F_n}u'_n|=\epsilon>0$.
Tetapi ini bertentangan dengan 246Yg.}
%246G

\spheader 246Yi Dalam 246Yg, tunjukkan bahwa $\nu$ bersifat aditif terhitung.
\Hint{Tetapkan $\mu=\sum_{n=0}^{\infty}a_n|\nu_n|$ untuk suatu barisan
$\sequencen{a_n}$ bilangan positif tegas yang sesuai. Untuk
setiap $n$ pilih turunan Radon-Nikod\'ym $f_n$ dari $\nu_n$ terhadap $\mu$.
Tunjukkan bahwa $\{f_n:n\in\Bbb N\}$ terintegralkan secara seragam, sehingga
$\nu$ kontinu sejati.} (Ini adalah {\bf teorema
Vitali-Hahn-Saks}.)
%246G

\spheader 246Yj Misalkan $(X,\Sigma,\mu)$ ruang ukur sembarang, dan
$A\subseteq L^1(\mu)$. Tunjukkan bahwa pernyataan-pernyataan berikut
ekuivalen: (i) $A$ terbatas dalam norma $\|\,\|_1$; (ii)
$\sup_{u\in A}|\int_Fu|<\infty$ untuk setiap $\mu$-atom $F\in\Sigma$ dan
$\limsup_{n\to\infty}\sup_{u\in A}|\int_{F_n}u|<\infty$ untuk setiap barisan
saling lepas $\sequencen{F_n}$ dari himpunan-himpunan terukur berukuran
berhingga; (iii) $\sup_{u\in
A}|\int_Eu|<\infty$ untuk setiap $E\in\Sigma$. \Hint{ tunjukkan bahwa
$\sequencen{a_nu_n}$ terintegralkan secara seragam bila
$\lim_{n\to\infty}a_n=0$ dalam $\Bbb R$ dan $\sequencen{u_n}$ suatu barisan
dalam $A$. }

%246G

\spheader 246Yk Misalkan $(X,\Sigma,\mu)$ adalah ruang ukur dan $A\subseteq
L^1(\mu)$ himpunan tak kosong. Tunjukkan bahwa pernyataan-pernyataan berikut
ekuivalen: (i) $A$ terintegralkan secara seragam; (ii) setiap kali
$B\subseteq
L^{\infty}(\mu)$ tidak kosong dan terarah ke bawah serta memiliki
infimum $0$ dalam $L^{\infty}(\mu)$ maka $\inf_{v\in B}\sup_{u\in A}|\int u\times
v|=0$. \Hint{Untuk (i)$\Rightarrow$(ii), catat bahwa $\inf_{v\in B}w\times
v=0$ untuk setiap $w\ge 0$ dalam $L^0$. Untuk (ii)$\Rightarrow$(i), gunakan
246G(iv).}
%246G

\spheader 246Yl Tetapkan $f(x)=e^{ix}$ untuk $x\in[-\pi,\pi]$. Tunjukkan bahwa
$|\int_Ef|\le 2$ untuk setiap $E\subseteq[-\pi,\pi]$.
%246K
}%end of exercises

\endnotes{\Notesheader{246} Sifat paling mencolok dari himpunan terintegralkan
secara seragam saya tunda sampai bagian berikutnya: himpunan-himpunan itu
relatif kompak secara lemah dalam $L^1$. Pembuktiannya memerlukan beberapa
gagasan yang tidak sederhana dari analisis fungsional dan topologi umum.
Bagian ini memuat hasil-hasil yang pada dasarnya berangkat dari teori ukuran.
Konsep, atau teknik, baru yang terpenting ialah `teorema barisan saling lepas'.
Contoh khas terdapat pada syarat (iii) dalam 246G, yang mengaitkan
keterintegralan seragam dengan perilaku fungsional pada barisan himpunan yang
saling lepas. Variasinya diberikan dalam 246Yd--246Yf, sedangkan
246Yg--246Yj memuat hasil-hasil lain yang dapat dibuktikan dengan metode
serupa. Hasil utama bagian berikutnya (247C) juga menggunakan barisan saling
lepas dalam pembuktiannya; barisan demikian akan muncul beberapa kali lagi
dalam Bab 35 pada volume berikutnya.

Frasa `terintegralkan secara seragam' semestinya berarti kira-kira `dapat
dihampiri secara seragam oleh fungsi sederhana'. Definisi 246A dapat diubah
ke bentuk semacam itu, tetapi menurut saya bentuk tersebut tidak terlalu
berguna. Sebaliknya, syarat (ii) dalam 246G kira-kira menyatakan sifat
`kontinu sejati secara seragam', jika anggota $L^1$ dipandang sebagai
fungsional kontinu sejati pada $\Sigma$, seperti dalam 242I. (Lihat 246Yi.)
Perhatikan bahwa dalam setiap pernyataan (ii)--(iv) dari 246G, atom-atom ukuran
harus ditangani secara khusus karena tidak dikendalikan oleh syarat utamanya.
Pada ruang ukur tak beratom, tentu terdapat penyederhanaan sebagaimana dalam
246Yd--246Yf.

Cara lain untuk membenarkan kata `seragam' dalam `terintegralkan secara
seragam' ialah dengan meninjau fungsional $\theta_w$, dengan $w\ge 0$ dalam
$L^1$, yang didefinisikan oleh $\theta_w(u)=\int(|u|-w)^+$ untuk $u\in L^1$.
Himpunan $A\subseteq L^1$ terintegralkan secara seragam jika dan hanya jika
$\theta_w\to 0$ secara seragam pada $A$ ketika $w$ meningkat dalam $L^1$
(246Xa). Jika sifat ini berlaku, kadang-kadang berguna bahwa sifat tersebut
dapat diperlihatkan dengan unsur $w$ yang dibangun langsung dari bahan yang
tersedia (lihat (iii) dalam 246Xa). Selanjutnya, himpunan
$A_{w\epsilon}=\{u:\theta_w(u)\le\epsilon\}$ selalu konveks, tertutup dalam
norma $\|\,\|_1$, dan `solid' (jika $u\in A_{w\epsilon}$ dan $|v|\le
|u|$, maka $v\in A_{w\epsilon}$) (246Cd). Himpunan ini tertutup terhadap
konvergensi titik demi titik dari barisan (246Ja), dan pada ruang ukur
semihingga juga tertutup dalam topologi konvergensi dalam ukuran (246Jd).
Pada ruang probabilitas, jika $w$ konstan, himpunan ini tertutup terhadap
ekspektasi bersyarat (246D) dan operator serupa (246Yc). Karena itu, kita dapat
mengharapkan bahwa setiap himpunan terintegralkan secara seragam termuat dalam
himpunan terintegralkan secara seragam yang tertutup terhadap beberapa jenis
operasi.

Sifat `seragam' lainnya tercantum dalam 246Yk. Norma $\|\,\|_{\infty}$ tidak
pernah kontinu terhadap urutan (dalam kasus-kasus yang menarik) sebagaimana
norma $\|\,\|_p$ lainnya (244Ye). Namun, subhimpunan terintegralkan secara
seragam dari $L^1$ menghasilkan seminorma-seminorma kontinu terhadap urutan
yang menarik pada $L^{\infty}$.

246J melengkapi hasil-hasil dari \S245. Dalam catatan untuk bagian itu saya
mengajukan pertanyaan berikut: jika $\sequencen{f_n}\to f$ hampir di mana-mana,
dengan cara apa $\sequencen{\int f_n}$ dapat gagal konvergen ke $\int f$?
Di sini diperoleh bahwa $\sequencen{\int|f_n-f|}\to 0$ jika dan hanya jika
$\{f_n:n\in\Bbb N\}$ terintegralkan secara seragam. Hal ini memperjelas
ungkapan `tidak ada bobot barisan yang hilang di tak hingga'. Secara umum,
untuk barisan, konvergensi dalam norma $\|\,\|_p$, dengan
$p\in\coint{1,\infty}$, adalah konvergensi dalam ukuran bagi barisan yang pangkat-$p$-nya
terintegralkan secara seragam (246Xh).




}%end of notes

\discrpage
